\section{Conclusions}

Empirical valence bond calculations show that the catalytic advantage of selenocysteine in GPX6 depends on the surrounding protein. Introducing Sec into mouse GPX6 lowers the proton-transfer barrier more than replacing Sec with Cys raises it in human GPX6. The same Sec/Cys exchange therefore has different effects in the two proteins.

The pathway calculations also reveal epistasis. Many routes connect human and mouse GPX6, and the triangular sequence projection does not indicate strict hysteresis. Nevertheless, some substitutions are consistently favored earlier than others, and their effects depend on whether position~49 contains Sec or Cys. GPX6 thus shows how changes elsewhere in a protein can compensate for replacement of a catalytic residue. The same approach can be applied to other Sec/Cys transitions and to enzyme design, where the effect of a catalytic substitution should be tested in more than one sequence context.
